\documentclass[a4paper, 12pt]{article}

\usepackage[utf8]{inputenc}
\usepackage[T1]{fontenc}
\usepackage{textcomp}
\usepackage{amssymb}
\usepackage{baskervald} \usepackage{newtxmath}
\usepackage{amsmath, amssymb}
\newtheorem{problem}{Problem}
\newtheorem{example}{Example}
\newtheorem{lemma}{Lemma}
\newtheorem{theorem}{Theorem}
\newtheorem{problem}{Problem}
\newtheorem{example}{Example} \newtheorem{definition}{Definition}
\newtheorem{lemma}{Lemma}
\newtheorem{theorem}{Theorem}
\DeclareMathAlphabet{\mathcal}{OMS}{cmsy}{m}{n}
\usepackage{parskip}

\begin{document}
    
\textit{Mr. D}. 

One can imaging Mr. D waking in the early morning, yet not too early, alone or
perhaps beside some sort of undemanding pet, a cat or perhaps even a tortoise. 
And one can \textit{only} imaging him waking up in this way. Hanging in some
dedicated section of his wardrobe, properly designated in relation to its
convenience and superiority over all other possible sections, he has prepared
the night before, with a complete lack of imagination, his outfit for
the day. His thoughts need not this form of foreseeing care, because he rarely
changes them. He is, our fin Mr. D, a truly harmless, even well-intentioned man, of
the special kind whose harmlessness seems mystically linked to a total lack of
competence, talent, or any kind of outstanding features. Looking at him, one
feels his soul must either be utterly shapeless, or a mere reflection of his
extraordinarily frail figure, and perhaps his small, timid eyes, his hookish and
protruding nose, his nervous smile themselves conform already something of a
spiritual visage, a kind of revelation. Nothing bad could be said about this
man: such a thing would already transgress against his indetermination. A fine
man, Mr. D, neither hated nor loved, seen and then forgotten, a child, a
well-intentioned soul.


\end{document}
